% Part: computability
% Chapter: recursive-functions
% Section: normal-form

\documentclass[../../../include/open-logic-section]{subfiles}

\begin{document}

\olfileid{cmp}{rec}{nft}
\olsection{نارمل روپ دا قضیہ}

\begin{thm}[کلینی دا نارمل روپ دا قضیہ (Kleene's Normal Form Theorem)]
\ollabel{thm:kleene-nf}
اک بنیادی بازگشتی تعلق $T(e, x, s)$ تے اک بنیادی
بازگشتی فنکشن $U(s)$ موجود نیں، جیہناں دی ایہہ
خاصیت اے: جے $f$ کوئی وی جزوی بازگشتی فنکشن
ہووے، تاں کسے~$e$ لئی،
\[
f(x) \simeq U(\umin{s}{T(e, x, s)})
\]
ہر $x$ لئی اے۔
\end{thm}

\begin{explain}
نارمل روپ دے قضیے دا ثبوت پیچیدہ اے، پر بنیادی
خیال سادہ اے۔ ہر جزوی بازگشتی فنکشن دا اک
\emph{اشاریہ (index)}~$e$ ہوندا اے؛ بدیہی طور
اُتے ایہہ اک عدد اے جیہڑا اوہدے پروگرام یا تعریف
دا کوڈ اے۔ جے $f(x)
\fdefined$، تاں حساب دا ریکارڈ باقاعدہ طریقے نال
رکھیا تے کسے عدد~$s$ وچ کوڈ کیتا جا سکدا اے؛
تے ایہہ گل کہ $s$،~$f$ دے اُس حساب دا کوڈ
اے جیہدا انپٹ~$x$ اے، صرف $x$ تے تعریف~$e$ ورت
کے بنیادی بازگشتی طریقے نال چیک کیتی جا سکدی
اے۔ ایس لئی تعلق~$T$، یعنی «اشاریہ~$e$ والے
فنکشن دا انپٹ~$x$ لئی اک حساب اے، تے $s$
اُس حساب دا کوڈ اے»، بنیادی بازگشتی اے۔ حساب
دا پورا ریکارڈ~$s$ دتا ہووے، تاں~$s$ دا
«نتیجہ»~$f(x)$ دی قیمت اے، تے اوہ وی~$s$
توں بنیادی بازگشتی طریقے نال حاصل کیتی جا سکدی اے۔

نارمل روپ دا قضیہ دسدا اے کہ کسے وی جزوی
بازگشتی فنکشن دی تعریف لئی صرف اک غیر محدود
تلاش دی لوڑ اے۔ بنیادی طور اُتے اسی سارے عدداں
وچ تلاش کر سکدے آں، جدوں تک اک ایہو جیہا عدد
نہ لبھے جیہڑا اشاریہ~$e$ والے فنکشن دے
انپٹ~$x$ لئی حساب دا کوڈ ہووے۔ اسی عدداں~$e$
نوں جزوی بازگشتی فنکشناں دے «ناں» دے طور اُتے
ورت سکدے آں، تے $\cfind{e}$ اوس فنکشن~$f$
لئی لکھ سکدے آں جیہڑا قضیے دی مساوات نال
تعریف ہوندا اے۔ یاد رکھو کہ کسے وی جزوی بازگشتی
فنکشن دے اک توں ودھ اشاریے ہو سکدے نیں—اصل
وچ، ہر جزوی بازگشتی فنکشن دے لامتناہی اشاریے نیں۔
\end{explain}
\end{document}
